% Part: lambda-calculus
% Chapter: introduction
% Section: fixed-point-combinator

\documentclass[../../../include/open-logic-section]{subfiles}

\begin{document}

\olfileid[fr]{lam}{int}{fix}
\olsection{Combinateurs de point fixe}

Supposons donné un lambda-terme~$g$ et cherchons
un autre terme~$k$ tel que $k$ soit
$\beta$-équivalent à~$gk$. Définissons
\[
\fn{diag}(x)=xx
\]
et
\[
l(x)=g(\fn{diag}(x))
\]
en choisissant $x\notin\FV{g}$. Autrement dit,
$l$ est le terme $\lambd[x][g(xx)]$. Prenons
$k=ll$. Nous obtenons
\begin{align*}
k & = (\lambd[x][g(xx)])(\lambd[x][g(xx)]) \\
& \red g((\lambd[x][g(xx)])(\lambd[x][g(xx)])) \\
& = gk.
\end{align*}
Si nous prenons
\[
Y=\lambd[g][((\lambd[x][g(xx)])(\lambd[x][g(xx)]))],
\]
alors $Yg$ et $g(Yg)$ se réduisent à un même
terme; ainsi $Yg\equiv_\beta g(Yg)$. C'est
le \emph{combinateur de Curry}. Si, au contraire,
nous prenons
\[
Y=(\lambd[xg][g(xxg)])(\lambd[xg][g(xxg)]),
\]
alors $Yg$ se réduit à $g(Yg)$, ce qui est
plus fort. Cette seconde version de~$Y$
est le \emph{combinateur de Turing}.

\end{document}
